% s-168-p168-m002; original PDF 168, 169, 170
\Aphorism{s-168-p168-m002}{257}
Wenn wir blind würden für die Gesichter der Menschen, taub für ihre Stimmen, und nur das Spiel ihrer Hände sähen, so würde unsere Welt im ersten Augenblick wohl gewaltig einbüssen an Mannigfaltigkeit und Ordnung. Aber auch dieses Chaos würde sich bald gliedern, und in kurzer Zeit, glaube ich, wären wir wieder soweit, nicht nur Berufsgattungen und gesellschaftliche Klassen, nicht nur Alter und Geschlecht, sondern geradezu Individuen unterscheiden zu können. Auch die Hand ist eine Physiognomie und hat eine, auch sie erfährt, denkt über das Erfahrene nach, widerspricht oder passt sich an, und nicht nur im Groben und Ungefähren, sodass man Arbeiter und Taschendiebe, Pfarrer und Soldaten, Herren und Bediente, Damen und Hausmütter an den Händen erkennt, sondern bis ins Feinste und Persönlichste — man muss nur Augen für dieses Spiel allerzartester Anzeichen haben. Die haben wir nicht, weil es uns mit den Gesichtern bequemer gemacht ist; in deren Ermanglung wären wir biologisch gezwungen, uns an die Hände zu halten, also auf Länge, Dicke und Bewegung der Finger, besonders des Daumens, auf Steifheit oder Lockerheit der Gelenke, auf die Zeichnung des Handtellers, auf Runzeln und Behaarung des Handrückens zu achten. Wir würden mit Sicherheit sagen können: diese Hand ist das Segnen, diese das Fechten, diese das Zugreifen gewöhnt; diese hat sich oft zur Faust geballt, diese liebt sich hohl zu machen, sei es um einen Trunk oder ein „Trinkgeld“ zu schöpfen; diese gehört Einem, der seiner Logik mit Fingerzeigen nachhelfen muss; hier widersprechen sich rechte und linke Hand, so coquett wie sich nur Mund und Augen eines Antlitzes widersprechen können. Von da bis zum Erkennen einer Person ist kein so grosser Schritt: man hat nur Alles, was die Hand einzeln verräth, wieder zum Gesammtbilde zu addiren und dies mit der Geschwindigkeit, mit der wir ein Antlitz als Ganzes sehen; so ist die Chirognomik als Praxis vollendet, wie sie es als Theorie darum noch lange nicht zu sein braucht. In Vorahnung dieses Zustandes trägt die gute Gesellschaft Handschuhe, d. h. man maskirt die persönliche Hand, die Indiscretionen begehen könnte, zu einer nichtssagenden uniformen Allerweltshand; von der ganzen reichen Symptomatik bleibt nur eine einzige Nummer übrig. Und eben darum muss Jedermann schreiben können, damit die Unterschiede in der Morphologie und Anatomie der Hände sich ausgleichen; Taschendieb und Arbeiter, Kriegsmann und Priester, Alle sollen eine mittlere Schreibhand haben.
\par

